%--------------------------
\begin{frame}{Question 5}

\begin{block}{}
Un nouvel État contractant :

\begin{itemize}
  \item[$\bullet$] a. n’est inclus dans la désignation globale et automatique contenue dans la requête 
  que s’il est devenu lié par le PCT avant la date à laquelle la demande internationale 
  est déposée. 
  \item[$\circ$] b. peut être élu dans la demande d’examen préliminaire international s’il est devenu lié 
  par le PCT à la date à laquelle la demande d’examen préliminaire a été déposée. 
  \item[$\circ$] c. peut être valablement inclus dans une désignation pour l’ouverture d’une phase 
  nationale s’il est devenu lié par le PCT à la date à laquelle la phase nationale est 
  engagée. 
  \item[$\circ$] d. Aucune des réponses a, b ou c n’est correcte.
\end{itemize}

\end{block}

\end{frame}

\begin{frame}{Question 5}

\begin{itemize}
  \item[$\bullet$] \textbf{a.} \; Correct : la désignation automatique ne couvre que les États
  qui sont liés par le PCT à la date de dépôt international
  (\pctrule{4}.9.a.i)).

  \item[$\circ$] \textbf{b.} \; Faux : un État ne peut être élu que s’il était déjà désigné,
  donc lié par le PCT à la date du dépôt international, pas seulement à la date de la demande
  Chapitre II (\pctart{31}.4.a)).

  \item[$\circ$] \textbf{c.} \; Faux : un État devenu contractant après la date de dépôt
  n’est valablement inclus en phase nationale que s'il était désigné donc lié
  par le PCT à la date de dépôt international (\pctart{22}.1).

  \item[$\circ$] \textbf{d.} \; Faux puisque la réponse \textbf{a} est correcte.
\end{itemize}

\end{frame}
